\subsection{GPT-3、The Pile、Gopher、LLaMA：数据配方开始变成公开工程}

GPT-3 使用 Common Crawl、WebText2、Books1/Books2 和 Wikipedia，约 400B tokens。其 Common Crawl 处理方式值得注意：训练了一个质量分类器，用 WebText、Wikipedia、Books1、Books2 作为正例，将其余 Common Crawl 内容作为负例，然后用这个分类器对整个 Common Crawl 做打分过滤。此外还做了模糊去重（包括与 WebText 和评测基准的去重）。

The Pile 则是开源社区对 GPT-3 封闭配方的回应，由大量志愿者在 Discord 上协调完成，汇聚 22 个高质量域，共 825GB 文本（约 275B tokens）。它包括 Pile-CC（使用 WARC + jusText 转换，优于 WET）、PubMed Central（500 万篇 NIH 资助的公开论文）、arXiv（使用 LaTeX 源码）、甚至安然公司邮件（500 名高管的 50 万封邮件，2002 年调查期间公开）。

Gopher 的 MassiveText 使用了 MassiveWeb（自建网页过滤）、C4、书籍、新闻、GitHub 和 Wikipedia，总计 10.5TB 文本，但 Gopher 实际只训练了 300B tokens（约 12\%）。MassiveWeb 的过滤使用手工规则而非分类器（例如”80\% 的词至少包含一个字母字符”），并使用 Google SafeSearch 做毒性过滤（而非词表）。

LLaMA 的数据混合又进一步把”多源拼接 + 质量过滤 + 近重复删除”做成了工程常态，使用 CCNet 处理的 Common Crawl、C4、GitHub（只保留宽松许可证）、Wikipedia（20 种语言）、Project Gutenberg 和 Books3（来自 The Pile）、arXiv、Stack Exchange（28 个最大站点），共 1.2T tokens。

\subsubsection{主要预训练数据集对比}

\begin{table}[H]
\centering
\small
\begin{tabular}{p{0.12\textwidth}p{0.07\textwidth}p{0.18\textwidth}p{0.1\textwidth}p{0.22\textwidth}p{0.12\textwidth}}
\toprule
\textbf{数据集} & \textbf{年份} & \textbf{主要来源} & \textbf{Tokens} & \textbf{使用的模型} & \textbf{是否公开} \\
\midrule
WebText & 2019 & Reddit 外链网页 & $\sim$10B & GPT-2 & 否 \\
C4 & 2019 & CC 单快照 + 规则 & 156B & T5 & 是 \\
CCNet & 2019 & CC + Wiki 过滤 & 可变 & XLM-R, BERT & 是 \\
GPT-3 data & 2020 & CC + books + wiki & 400B & GPT-3 & 否 \\
The Pile & 2021 & 22 域混合 & 275B & GPT-J, NeoX & 是 \\
MassiveText & 2021 & 多源 + 规则过滤 & 10.5TB & Gopher & 否 \\
LLaMA data & 2023 & CCNet+C4+GitHub 等 & 1.2T & LLaMA & 配方公开 \\
RedPajama v1 & 2023 & LLaMA 配方复现 & 1.2T & -- & 是 \\
RefinedWeb & 2023 & CC + trafilatura & 5T & Falcon & 部分(600B) \\
FineWeb & 2024 & 95 个 CC 快照 & 15T & -- & 是 \\
Dolma & 2024 & 多源(Reddit 等) & 3T & OLMo & 是 \\
DCLM & 2024 & CC + 分类器过滤 & 3.8T & -- & 是 \\
Nemotron-CC & 2024 & CC + 重写 + 集成 & 6.3T & Nemotron & 部分 \\
\bottomrule
\end{tabular}
\caption{主要预训练数据集对比：从 10B tokens 到 15T tokens，规模增长了三个数量级}
\end{table}

\subsubsection{数据过滤方法对比}

\begin{table}[H]
\centering
\begin{tabular}{p{0.15\textwidth}p{0.2\textwidth}p{0.23\textwidth}p{0.15\textwidth}p{0.15\textwidth}}
\toprule
\textbf{过滤方法} & \textbf{原理} & \textbf{优点} & \textbf{缺点} & \textbf{代表} \\
\midrule
规则过滤 & 手写规则（长度、标点、脏词等） & 透明、可解释、无偏差引入 & 误杀/漏网多，难以覆盖所有情况 & C4, Gopher \\
Perplexity 过滤 & 用 Wikipedia 训练 LM，低 PPL 即高质量 & 有统计基础，自动化程度高 & 偏向 Wikipedia 风格，排斥口语和创意 & CCNet \\
分类器过滤 & 训练二分类器区分”好/坏”文本 & 灵活，可定制正负例 & 正例选择决定偏见方向 & DCLM, GPT-3 \\
集成过滤 & 多个分类器/打分器投票 & 降低单一偏见 & 计算成本高 & Nemotron-CC \\
\bottomrule
\end{tabular}
\caption{四种主要过滤方法各有适用场景，现代系统通常组合使用}
\end{table}

\begin{knowledgebox}{为什么公开配方有价值}
它让研究者能把”模型效果”拆成”数据选择”的结果，而不是只看最终指标。数据公开得越清楚，越能比较不同 filtering / mixing 策略的真实影响。值得注意的是，最新一代模型的训练 token 量已经远超上表：Llama 3 训练了 15T tokens，Qwen3 训练了 36T tokens。数据规模的增长速度甚至超过了模型参数的增长。
\end{knowledgebox}

